\section{Solutions for Chapter~\ref{ch:related-categories}}

\paragraph{Exercise~\ref{ex:noncyclic-factorization}.}
\phantomsection
\label{sol:ex:noncyclic-factorization}
Choose nonzero $h\in H$ and $k\in K$. The reciprocal conditions
on $(h,0)$ and $(0,k)$ read
\[
 U(h)U(-h)V(0)^2=\delta,\qquad
 U(0)^2V(k)V(-k)=\delta.
\]
Since $U(0)^2V(0)^2=W(0,0)^2=1$, multiplying gives
$U(h)U(-h)V(k)V(-k)=\delta^2$. The condition at $(h,k)$
says that the same product equals $\delta$. Nonvanishing forces
$\delta=1$. But $\delta^2=n\delta+1$ with $n>0$ does not
permit $\delta=1$. Coupling the coordinates or changing the
origin structure is therefore necessary; the direct product of
independent reciprocal scalar functions cannot solve this ansatz.



\paragraph{Exercise~\ref{ex:noncyclic-eight-branch}.}
\phantomsection
\label{sol:ex:noncyclic-eight-branch}
Put $s=t^2$. Since $s>0$ and $s^2+8s-1=0$,
$s=\sqrt{17}-4$ lies strictly between zero and one.
Dividing the second eliminant by $2tr^2$ gives
$r^{-2}-r^2=(3-s)/(2t)>0$, hence $0<r<1$. Squaring and
adding four yields
\[
 (r^2+r^{-2})^2
 =\frac{(3-s)^2+16s}{4s}
 =\frac{9+10s+s^2}{4s}
 =\left(\frac{9+s}{2}\right)^2,
\]
where $s^2+8s-1=0$ reduces the last equality. Since
$r^2+r^{-2}>0$ and $9+s>0$, the positive square root gives
$r^2+r^{-2}=(9+s)/2$. Thus
\[
 (r+r^{-1})^2=(13+s)/2=((5+s)/2)^2,
\]
and positivity again gives $r+r^{-1}=(5+s)/2$.
Solving the resulting quadratic and choosing the root below one
gives
\[
 r=\frac{5+s-\sqrt{(5+s)^2-16}}4
  =\frac{1+\sqrt{17}-\sqrt{2+2\sqrt{17}}}{4}
  =\lambda_8^2.
\]
For completeness, the supplied identity $tr^3+2r=1$ can also
be derived from this positive branch. The relation $s^2+8s=1$
gives
\[
 (5+s)^2-16=s(9+s)^2.
\]
Since $t>0$, the selected radical is $t(9+s)$, and hence
\[
 r=\frac{5+s-t(9+s)}4
   =\frac{t(5+s)+2}{t(11+s)+4}.
\]
The second equality follows by cross-multiplication and reduction
using $t^2=s$ and $s^2+8s=1$; its denominator is positive.
Writing $C=(5+s)/2$, the quadratic $r^2-Cr+1=0$ gives
\[
 r^3=(C^2-1)r-C
     =\frac{11+s}{2}r-\frac{5+s}{2}.
\]
Consequently
\[
 tr^3+2r
 =\frac{t(11+s)+4}{2}r-\frac{t(5+s)}2=1.
\]
Multiplication of the first eliminant by
$r$ now gives $y^2r^3=tr^3+2r=1$. The positive sample $y$ is
therefore $r^{-3/2}=\lambda_8^{-3}$. The signs of $r,y$ and
the positive square-root choices are analytic branch information
that the polynomial equations alone do not supply.



\paragraph{Exercise~\ref{ex:noncyclic-smith}.}
\phantomsection
\label{sol:ex:noncyclic-smith}
At $m=3$, the relations are $D-1=4+2\omega$ and
$\omega^2=\omega+9$. Multiplication of the basis vectors
$1,\omega$ gives $4+2\omega$ and $18+6\omega$,
so the matrix is
\[
 \begin{pmatrix}4&18\\2&6\end{pmatrix}.
\]
Its determinant is $24-36=-12$ and the gcd of its entries is
$2$. For a full-rank two-by-two integer matrix, the first Smith
invariant is that gcd and the product of the two invariants is
the absolute determinant. They are therefore $2,6$. The
quotient has $12$ elements, but every element has order dividing
six. The cyclic group of order twelve has an element of order
twelve, so the cardinality alone loses essential group data.



\paragraph{Exercise~\ref{ex:EH-moment-projections}.}
\phantomsection
\label{sol:ex:EH-moment-projections}
Let $Y=S^2-f$. The support condition gives
\[
 \operatorname{tr}(Y^*Y)
 =\operatorname{tr}(S^4)-2\operatorname{tr}(S^2)
   +\operatorname{tr}(f)=0.
\]
Faithful positivity implies $Y=0$. Expanding with $S^2=f$
and $fS=S$ gives
\[
 \left(\frac{f\pm S}{2}\right)^2
 =\frac{2f\pm2S}{4}=\frac{f\pm S}{2},\qquad
 \frac{f+S}{2}\frac{f-S}{2}=0.
\]
They are self-adjoint and sum to $f$, and their traces are
$\operatorname{tr}(f)/2$ because $\operatorname{tr}(S)=0$.
For an indefinite pairing a nonzero $Y$ may have zero squared
pairing, so the first inference to $Y=0$ fails.
For example, on $\C\oplus\C$ take
$\operatorname{tr}(a,b)=a-b$, $f=(1,1)$ and $S=0$.
The stated trace equalities and support conditions hold, but
$S^2\ne f$. The trace pairing is nondegenerate and indefinite;
it cannot turn a zero squared pairing into a zero vector.



\paragraph{Exercise~\ref{ex:EH-word-tensor}.}
\phantomsection
\label{sol:ex:EH-word-tensor}
If $L^kw=w$ with odd $k$, the rotation rule gives
$p_w=p_{L^kw}=(-1)^kp_w=-p_w$, hence $p_w=0$ over $\C$.
Every letter $2$ may be replaced by the sum of letters $0$ and
$1$ with a minus sign, decreasing the number of $2$'s. This
determines all coefficients from binary ones. Equivalently,
the uniquely specified multilinear tensor evaluated at
$u_2=-u_0-u_1$ performs every replacement at once; independence
of the replacement order follows from multilinearity. In the
given example,
\[
 p_{00000002}=-p_{00000000}-p_{00000001}
 =-c_1=2\lambda^4+\lambda^2-9.
\]



\paragraph{Exercise~\ref{ex:EH-reflection-F}.}
\phantomsection
\label{sol:ex:EH-reflection-F}
The defining cubic is negative at $x=4$ and tends to $+\infty$,
so its selected largest root satisfies $x>4$. Consequently
$a=x-2>2$, $y=a+1>0$, and $z=a^2+a-1>0$; the displayed
square roots make $v$ a real vector.
With $y=a+1$ and $z=a^2+a-1$,
\[
 v^Tv=\frac{a^2+y+z}{2ay}=1.
\]
The two requested entries are
\[
 -2v_1v_2=-\frac1{\sqrt y},\qquad
 -2v_2v_3=\frac{\sqrt{z/y}}a.
\]
The other entries follow in the same way. Since $v^Tv=1$,
$(I-2vv^T)^2=I-4vv^T+4v(v^Tv)v^T=I$.
The matrix is real symmetric, hence orthogonal and unitary.
Its eigenvalues are $-1,1,1$, so its determinant is $-1$.
This proves a matrix identity. Identifying it as a particular
category's associator additionally requires the specified
trivalent bases and their planar evaluations; many unrelated
orthogonal matrices satisfy the same unitarity equation.
